% page 47 of the facsimile (image p-046 of the scan)
% block 167.6 x 242.0 mm ; left margin 34.4 mm
\pgc{12.700mm}{0.000mm}{
\l{\cel{54}39}
\l{}
\l{\sou{aristokrata}. Qua apartenas a ta formo di guvernado en olqua}
\l{\cel{3}la autoritato guvernala esas la privilejo di la nobeli.-DEFIRS.}
\l{}
\l{\sou{aristolokio}. (\sou{bot.}) Planto lignoza di qua la radiko uzesas kom}
\l{\cel{3}tonizivo ed apertivo. - L.\cel{1}\sou{aristolochia}.\cel{2}DEFIS.}
\l{}
\l{\sou{aritmetiko}. (\sou{matem.}) Cienco di la nombri, di kalkulado. - DEFIRS.}
\l{}
\l{\sou{aritmografo}. Aparato por kalkular mekanike ; kalkulo-mashino.}
\l{\cel{3}- DEFIRS.}
\l{}
\l{\sou{aritmologio}. Nomo di la cienco generala di la nombri, di la}
\l{\cel{3}mezuro di la +grandori, irgequalan esas oli. - DEFIRS.}
\l{}
\l{\sou{aritmometro}. Kalkulo-linealo. - DEFIRS.}
\l{}
\l{\sou{arivar}. (\sou{netrans.,an}). Abutar an la fino-punto di sua iro.}
\l{\cel{3}- EFIS.}
\l{}
\l{\sou{arjento}.(\sou{kemio}). Metalo blanka, briloza, quan onu uzas por}
\l{\cel{3}fabrikar moneto-peci, vazi, juveli, adjuntante ad olu}
\l{\cel{3}aloyuro de kupro. -\cel{1}\sou{Simbolo kemiala}\cel{1}:\cel{2}Ag.\cel{2}- EFIS.}
\l{}
\l{\sou{arjentano}.\cel{2}Aloyuro de nikelo kun kupro o zinko. - DEFIS.}
\l{}
\l{\sou{arko}. I. Parto limitoza di kurta lineo irgaspeca. - II.}
\l{\cel{3}Konstrukturo qua havas la formo di kurvo. - III. Lumo}
\l{\cel{3}dazlanta qua spricas inter du peci de karbono ligita}
\l{\cel{3}an la du elektrodi di generatoro elektrala. - DEFIRS.}
\l{}
\l{\sou{arkabuzo}. Olima paf-armo portebla. - DEFIRS.}
\l{}
\l{\sou{arkado}. Konstrukturo qua havas la formo di arko, qua su}
\l{\cel{3}apogas a pilastri o koloni qui posibligas subpaso. -}
\l{\cel{3}DEFIRS.}
\l{}
\l{\sou{arkaika}. Qua evas de la komenco di arto. - DEFIRS.}
\l{}
\l{\sou{arkaismo}. Expresuro, vort-aranjo qua esas apene uzata pro}
\l{\cel{3}olua quaza antiqueso. - DEFIRS.}
\l{}
\l{\sou{arkano}. Sekreto ciencala (o pseudo-ciencala) quan onu re-}
\l{\cel{3}velas o docas nur ad iniciati. (Ex.: la arkani di al-}
\l{\cel{3}kemio). - DEFIRS.}
\l{}
\l{\sou{arkegonio.}\cel{2}(\sou{bot.}) Genit-organo feminala di la planti hepa-}
\l{\cel{3}tikala, di la mosi e di la kriptogami vaskulala e gimno-}
\l{\cel{3}sperma. - DEFIS.}
\l{}
\l{\sou{arkeologio}. Cienco pri la kozi antiqua. - DEFIRS.}
\l{}
\l{\sou{aritmomanciar}. (\sou{netrans.})\cel{2}Asertar savar (per procedi super}
\l{\cel{2}natura) to quo esas celita en la pasinto, la prezento e la}
\l{\cel{2}futuro - per nombri. - DEFIRS.}
}
